\begin{lemma}[Existence of the $L^1$ extension satisfying the weak bound]
There is a linear map
\[
  T : L^1(\gamma_n;\mathbb{R}) \longrightarrow
      \mathcal{M}(\gamma_n;\mathbb{R}^n),
\]
where $\mathcal{M}(\gamma_n;\mathbb{R}^n)$ denotes the space of
$\gamma_n$-almost-everywhere equivalence classes of measurable
$\mathbb{R}^n$-valued functions, such that
\begin{enumerate}
\item $T$ agrees with $R_\gamma$ on $L^2(\gamma_n)$: for every
$v \in L^2(\gamma_n)$ one has $T(\iota v) = R_\gamma v$ almost everywhere, with
$\iota$ the inclusion of \noderef{ToL1} and $R_\gamma$ as in
\noderef{RieszTransform}; and
\item $T$ satisfies the weak bound
\[
  t\,\gamma_n\bigl(\{x : |Tf(x)| > t\}\bigr) \leq 2\,\|f\|_{L^1(\gamma_n)}
  \qquad\text{for every } f \in L^1(\gamma_n) \text{ and } t > 0 .
\]
\end{enumerate}
\end{lemma}
\begin{proof}
\emph{Step 0: a Cauchy-in-measure estimate.}
Let $v, w \in L^2(\gamma_n)$ and $\varepsilon > 0$.  By \noderef{RieszTransform}
the map $R_\gamma$ is linear on $L^2(\gamma_n)$, so
$R_\gamma v - R_\gamma w = R_\gamma(v-w)$, and \noderef{MainBound} applied to the
real $L^2$ function $v - w$ at the level $\varepsilon$ gives
\begin{equation*}
  \gamma_n\bigl(\{|R_\gamma v - R_\gamma w| > \varepsilon\}\bigr)
    \leq \frac{2}{\varepsilon}\,\|\iota(v-w)\|_{L^1(\gamma_n)} . \tag{1}
\end{equation*}

\emph{Step 1: truncations approximate in $L^1$.}
Let $f \in L^1(\gamma_n)$ and put $f_k = \max\{-k, \min\{f,k\}\}$ for
$k \in \mathbb{N}$.  Each $f_k$ is measurable with $|f_k| \leq k$, hence bounded,
hence in $L^2(\gamma_n)$ because $\gamma_n$ is a probability measure
(\noderef{GaussIsProb}).  Moreover $|f_k| \leq |f|$ and $f_k \to f$ pointwise, so
$|f - f_k| \leq 2|f| \in L^1(\gamma_n)$ and $|f-f_k| \to 0$ pointwise; dominated
convergence gives $\|f - f_k\|_{L^1(\gamma_n)} \to 0$.  In particular
$(\iota f_k)_k$ is Cauchy in $L^1(\gamma_n)$.

\emph{Step 2: an almost-everywhere convergent subsequence.}
Using Step 1, choose indices $k_1 < k_2 < \cdots$ with
$\|f_{k_{j+1}} - f_{k_j}\|_{L^1(\gamma_n)} \leq 2^{-3j}$ for every $j \geq 1$
(possible since the sequence is Cauchy in $L^1$).  Applying (1) with
$v = f_{k_{j+1}}$, $w = f_{k_j}$ and $\varepsilon = 2^{-j}$,
\[
  \gamma_n\bigl(\{|R_\gamma f_{k_{j+1}} - R_\gamma f_{k_j}| > 2^{-j}\}\bigr)
    \leq 2\cdot 2^{j}\cdot 2^{-3j} = 2^{1-2j} .
\]
The right-hand side is summable in $j$, so by the Borel--Cantelli lemma the set of
$x$ lying in infinitely many of the sets
$E_j = \{|R_\gamma f_{k_{j+1}} - R_\gamma f_{k_j}| > 2^{-j}\}$ is
$\gamma_n$-null.  Hence for $\gamma_n$-a.e.\ $x$ there is $J(x)$ with
$|R_\gamma f_{k_{j+1}}(x) - R_\gamma f_{k_j}(x)| \leq 2^{-j}$ for all
$j \geq J(x)$; the tail sums $\sum_{j \geq J} 2^{-j}$ are finite, so
$(R_\gamma f_{k_j}(x))_j$ is a Cauchy sequence in $\mathbb{R}^n$ and converges to
a finite limit.  Define $Tf$ to be the $\gamma_n$-a.e.\ pointwise limit of
$R_\gamma f_{k_j}$ (set to $0$ on the null set where convergence fails); it is a
measurable $\mathbb{R}^n$-valued function, hence an element of
$\mathcal{M}(\gamma_n;\mathbb{R}^n)$.

\emph{Step 3: convergence in measure of the full sequence, and well-definedness.}
By Step 2 and Egorov-free elementary reasoning, $R_\gamma f_{k_j} \to Tf$
$\gamma_n$-a.e., hence in measure (on the finite measure space $\gamma_n$,
a.e.\ convergence implies convergence in measure).  Combined with (1) and Step 1,
which make $(R_\gamma f_k)_k$ Cauchy in measure, the full sequence
$R_\gamma f_k$ converges to $Tf$ in measure: given $\varepsilon > 0$,
\[
  \gamma_n(\{|R_\gamma f_k - Tf| > 2\varepsilon\})
    \leq \gamma_n(\{|R_\gamma f_k - R_\gamma f_{k_j}| > \varepsilon\})
       + \gamma_n(\{|R_\gamma f_{k_j} - Tf| > \varepsilon\}),
\]
and both terms can be made small, the first by (1) and Step 1, the second by
a.e.\ convergence.

If $(\tilde f_k) \subseteq L^2(\gamma_n)$ is any other sequence with
$\iota\tilde f_k \to f$ in $L^1(\gamma_n)$, then (1) applied to the pair
$f_k, \tilde f_k$ gives
$\gamma_n(\{|R_\gamma f_k - R_\gamma\tilde f_k| > \varepsilon\}) \leq
\tfrac{2}{\varepsilon}\|\iota f_k - \iota \tilde f_k\|_{L^1} \to 0$, so
$R_\gamma \tilde f_k \to Tf$ in measure as well, and a limit in measure is unique
up to a null set.  Hence $Tf$ does not depend on the approximating sequence.

\emph{Step 4: linearity.}
Let $f, h \in L^1(\gamma_n)$ and $\lambda \in \mathbb{R}$, and let
$(f_k), (h_k)$ be the truncation sequences of Step 1.  Then
$(f_k + h_k)$ and $(\lambda f_k)$ are sequences in $L^2(\gamma_n)$ whose
inclusions converge in $L^1(\gamma_n)$ to $f + h$ and $\lambda f$ respectively.
By linearity of $R_\gamma$ (\noderef{RieszTransform}),
$R_\gamma(f_k+h_k) = R_\gamma f_k + R_\gamma h_k$ and
$R_\gamma(\lambda f_k) = \lambda R_\gamma f_k$.  Convergence in measure is
compatible with addition and scalar multiplication, so the left sides converge in
measure to $T(f+h)$ and $T(\lambda f)$ by Step 3, while the right sides converge
in measure to $Tf + Th$ and $\lambda\,Tf$.  Uniqueness of limits in measure gives
$T(f+h) = Tf + Th$ and $T(\lambda f) = \lambda Tf$.

\emph{Step 5: agreement with $R_\gamma$ on $L^2$.}
Let $v \in L^2(\gamma_n)$.  The constant sequence $\tilde f_k = v$ lies in
$L^2(\gamma_n)$ and satisfies $\iota \tilde f_k \to \iota v$ in $L^1(\gamma_n)$
trivially.  By the independence established in Step 3,
$T(\iota v) = \lim_k R_\gamma \tilde f_k = R_\gamma v$ $\gamma_n$-a.e.

\emph{Step 6: the weak bound.}
Fix $f \in L^1(\gamma_n)$ and $t > 0$, and let $(f_{k_j})$ be the subsequence of
Step 2, so $R_\gamma f_{k_j} \to Tf$ $\gamma_n$-a.e.  At a point $x$ with
$|Tf(x)| > t$ there is a strict margin $|Tf(x)| - t > 0$, and
$|R_\gamma f_{k_j}(x)| \to |Tf(x)|$, so $|R_\gamma f_{k_j}(x)| > t$ for all large
$j$.  Hence $\gamma_n$-a.e.
\[
  \mathbf{1}_{\{|Tf|>t\}} \leq \liminf_{j\to\infty}\mathbf{1}_{\{|R_\gamma f_{k_j}|>t\}} .
\]
Fatou's lemma, applied to these nonnegative measurable functions, gives
\[
  \gamma_n(\{|Tf|>t\})
    \leq \liminf_{j} \gamma_n\bigl(\{|R_\gamma f_{k_j}| > t\}\bigr) .
\]
By \noderef{MainBound} applied to $f_{k_j}$ at level $t$, each term is at most
$\tfrac{2}{t}\|\iota f_{k_j}\|_{L^1(\gamma_n)}$.  By Step 1,
$\|\iota f_{k_j}\|_{L^1(\gamma_n)} \to \|f\|_{L^1(\gamma_n)}$, so
\[
  \gamma_n(\{|Tf|>t\}) \leq \frac{2}{t}\,\|f\|_{L^1(\gamma_n)} .
\]
Multiplying by $t > 0$ gives the asserted bound.
\end{proof}
